% SPDX-License-Identifier: Apache-2.0
\section{\code{txhdl::comp}}
\label{sec:r-comp}

The composition module: everything a unit is made of, and the executor
that runs it. In order of appearance: clocks with their period and
phase, and the edge wait a process makes; a wire and a channel, each
split into a driving end that is not \code{Clone} and a reading end
that is, the channel's reading end with a \code{wait}; the
\code{Member} trait that lets \code{interface!} split a member without
knowing which it is; registers and memories, read
plainly and driven with a commit at the end of the step; \code{Regs},
\code{N} registers of one type as one field, since a unit is
\code{Default} and an array of a generic length has none, its
registers \code{name\_0} onward in the netlist and the trace through
\code{trace\_as} and \code{Port::entries} (issue 594); \code{Units}, \code{N}
child units of one type as one field, for the same reason, and
\code{chans}, \code{signals} and \code{Ends}, which hand out \code{N}
ends once each by index, so that a join of those children can give
child \code{i} the ends at \code{i} or at a neighbour's (issue 635); \code{mux};
the \code{Unit} trait, whose \code{run} takes the unit's inputs and
outputs as type parameters; configurations and the \code{config!}
macro; \code{join2} and \code{join\_all}, which give each child its
own process identity; \code{parallel!}, which keeps its futures in one
process and returns their outputs together; the executor with its
waits, \code{edge}, \code{until} and \code{tick}, whose rules
Section~\ref{sec:r-time} states; and \code{trace}, the probe registry
and the VCD and FST writers.

\lstinputlisting[caption={\code{lib/src/comp.rs}.},label={lst:r-comp}]{lib/src/comp.rs}
